\chapter{Phân tích, thiết kế và hiện thực}

Chương này gắn yêu cầu vào hiện thực. Mỗi khối thiết kế (cắt Điều, bảy nút,
FTS, JWT, lịch, Docker) được đọc theo cùng bốn câu như Chương~2: khối đó
\textbf{là gì}, \textbf{vì sao} chọn, \textbf{gắn vào hệ thống thế nào}, và
\textbf{chạy khi nào}.

\section{Tổng quan chức năng từ góc nhìn người dùng}

Người dùng tạo tài khoản, khai tên công ty, mã số thuế, số lao động, kỳ kê
khai GTGT. Bốn việc chính sau đó:

\textbf{Hỏi.} Gõ câu hỏi, thấy từng bước (đang tìm, đang lọc, đang viết), thấy
từng chữ hiện ra, thấy danh sách Điều bấm được. Hội thoại cũ nằm cột trái, sống
sót khi restart máy chủ.

\textbf{Đọc luật.} Lọc theo lĩnh vực, vào một văn bản, đi từ Chương xuống Điều,
bấm viện dẫn chéo, hoặc bấm ``Hỏi AI về Điều này'' để câu hỏi được điền sẵn
kèm số hiệu.

\textbf{Đưa hợp đồng.} Kéo file PDF/DOCX/TXT, tối đa 10~MB. Máy trích chữ ngay
(lỗi file lộ lúc tải, không đợi). Bấm soát xét thì chạy nền; trang hỏi trạng
thái mỗi vài giây. Kết quả là bảng: điều khoản, mức cao/trung bình/thấp, đề
xuất sửa, căn cứ.

\textbf{Nhìn lịch.} Các mốc trong năm, nhóm theo tháng, đánh dấu xong hoặc bỏ
qua. Sinh lại khi đổi số lao động hay đổi kỳ GTGT.

Người quản trị (tài khoản đăng ký đầu tiên khi bảng user còn trống được gắn vai
trò này) vào được kho văn bản, danh sách user, và trang chạy hàng loạt câu hỏi
để kiểm tra pipeline.

\section{Yêu cầu}

\subsection{Chức năng}

\begin{table}[htp]
\centering
\caption{Yêu cầu chức năng của hệ thống.}
\label{tab:yc-cn}
\begin{tabular}{p{1.6cm}p{5.2cm}p{7.0cm}}
\toprule
\textbf{Mã} & \textbf{Việc} & \textbf{Ghi chú khi cài} \\
\midrule
YC-01 & Đăng ký, đăng nhập, làm mới token, đổi mật khẩu & Access 15 phút,
refresh 7 ngày, hai loại token không dùng lẫn \\
YC-02 & Hỏi đáp có trích dẫn và stream & Ghi câu hỏi vào CSDL trước khi gọi
pipeline \\
YC-03 & Lịch sử hội thoại & Đổi tên, lưu trữ, xóa; không nhìn hội thoại người
khác (trả 404) \\
YC-04 & Danh mục và cây văn bản & Số hiệu có dấu \texttt{/} nên định tuyến
phải kiểu path \\
YC-05 & Tìm toàn văn không dấu & Ba tầng: số Điều, \texttt{ts\_rank}, trigram \\
YC-06 & Tải và soát xét hợp đồng & Trích chữ lúc tải; chấm điểm lúc soát xét \\
YC-07 & Lịch tuân thủ & Idempotent: gọi sinh lại không xóa việc đã đánh dấu xong \\
YC-08 & Quản trị corpus và user & Nhập JSON, reindex, khóa tài khoản \\
YC-09 & Chạy hàng loạt câu hỏi & Chỉ admin; JSON danh sách câu, không dùng hội thoại \\
\bottomrule
\end{tabular}
\end{table}

\subsection{Phi chức năng}

Thời gian: trang tra cứu phải trả trong khoảng dưới một giây trên máy phát triển
dùng cho đồ án; hỏi
đáp chấp nhận hàng chục giây nếu hiện tiến độ. Bảo mật: mật khẩu argon2, CORS
không phải \texttt{*}, giới hạn số lần gọi chat. Sẵn sàng: thiếu embedding thì
phần còn lại vẫn sống.

\begin{table}[htp]
\centering
\caption{Ngưỡng phi chức năng của hệ thống.}
\label{tab:yc-phi-cn}
\begin{tabular}{p{3.2cm}p{5.0cm}p{5.5cm}}
\toprule
\textbf{Nhóm} & \textbf{Ngưỡng} & \textbf{Cách kiểm} \\
\midrule
Độ trễ tra cứu & $< 1$ giây trên máy thử nghiệm với FTS & Gọi \texttt{/laws/search} \\
Độ trễ hỏi đáp & Chấp nhận hàng chục giây nếu có SSE tiến độ & Quan sát UI \\
Bảo mật mật khẩu & argon2id, tối thiểu 8 ký tự & Đăng ký trên UI \\
Cách ly dữ liệu & Hội thoại theo user; hợp đồng/lịch theo org & 404 khi xem chéo \\
Sẵn sàng suy giảm & Thiếu embedding: health 200, chat 503 & Chạy Docker thiếu khóa \\
Giới hạn tần suất & Chat 20/phút; upload 30/giờ & slowapi \\
\bottomrule
\end{tabular}
\end{table}

\section{Kiến trúc ba lớp}

\begin{figure}[htp]
\centering
\hinhKienTruc
\caption{Ba lớp triển khai: trình duyệt $\rightarrow$ nginx $\rightarrow$ FastAPI;
PostgreSQL, Chroma và Gemini nằm phía sau API.}
\label{fig:kien-truc-he-thong}
\end{figure}

Lúc cài bằng Docker, trình duyệt chỉ nói chuyện với cổng frontend. Nginx vừa
đưa file tĩnh vừa chuyển \texttt{/api} sang backend, nên không có chuyện CORS
trên môi trường đó. Khi phát triển, Vite thực hiện chuyển tiếp tương tự trên
cổng 5173.

Logic RAG không nằm trong file router. Router chỉ nhận request, ghi hội thoại,
gọi agent. Cùng một hàm truy hồi được soát xét hợp đồng gọi lại, tránh nhân bản
prompt giữa hai module.

BM25 không để Postgres. Khởi động backend: đọc Điều từ Postgres, nếu file cache
khớp dấu vân tay tokenizer thì nạp cache, không thì tách từ lại và ghi cache.
Vector: nếu manifest (model, số chiều, số record) lệch thì nhúng lại, có chia
lô và nghỉ dần khi bị 429.

\section{Use case}

Use case mô tả tương tác giữa tác nhân và hệ thống để đạt một mục tiêu~\cite{thanh2024}.
Mỗi use case có tiền điều kiện, luồng chính, nhánh và lỗi. Đồ án dùng use case
làm cầu giữa yêu cầu ngôn ngữ tự nhiên và API.

\begin{table}[htp]
\centering
\caption{Tác nhân.}
\label{tab:vai-tro}
\begin{tabular}{p{3.0cm}p{5.4cm}p{5.4cm}}
\toprule
\textbf{Tác nhân} & \textbf{Người nào} & \textbf{Được làm} \\
\midrule
Chủ doanh nghiệp & Người mở tài khoản công ty & Cả bốn nghiệp vụ trong công ty
mình \\
Kế toán, Nhân sự & Người được mời vào cùng công ty & Cùng bốn nghiệp vụ; không
vào trang quản trị \\
Quản trị & Tài khoản đầu, hoặc được đổi role & Thêm quyền corpus, user, chạy
hàng loạt \\
Khách & Chưa đăng nhập & Chỉ đăng ký / đăng nhập \\
\bottomrule
\end{tabular}
\end{table}

\begin{figure}[htp]
\centering
\resizebox{0.95\linewidth}{!}{\hinhUseCase}
\caption{Sơ đồ use case. Chủ / kế toán / nhân sự dùng bốn nghiệp vụ; quản trị
thêm corpus và chạy hàng loạt câu hỏi.}
\label{fig:so-do-use-case}
\end{figure}

\begin{table}[htp]
\centering
\caption{Danh sách use case.}
\label{tab:danh-sach-uc}
\begin{tabular}{p{1.5cm}p{6.0cm}p{6.3cm}}
\toprule
\textbf{Mã} & \textbf{Tên} & \textbf{Ai} \\
\midrule
UC-01 & Tạo tài khoản kèm hồ sơ công ty & Khách \\
UC-02 & Đăng nhập, làm mới phiên & Mọi vai trò đã có tài khoản \\
UC-03 & Hỏi đáp pháp luật & Chủ, kế toán, nhân sự \\
UC-04 & Xem và đổi tên hội thoại & Người sở hữu hội thoại \\
UC-05 & Đọc văn bản theo cây Chương--Điều & Chủ, kế toán, nhân sự \\
UC-06 & Tìm chữ trong corpus & Chủ, kế toán, nhân sự \\
UC-07 & Tải hợp đồng và chạy soát xét & Người thuộc một công ty \\
UC-08 & Đọc kết quả soát xét & Người cùng công ty \\
UC-09 & Xem lịch, đánh dấu xong & Người thuộc một công ty \\
UC-10 & Nhập thêm văn bản, dựng lại chỉ mục & Quản trị \\
UC-11 & Khóa / đổi vai trò user & Quản trị \\
UC-12 & Chạy hàng loạt câu hỏi & Quản trị \\
\bottomrule
\end{tabular}
\end{table}

\begin{table}[htp]
\centering
\caption{Đặc tả UC-03 --- Hỏi đáp.}
\label{tab:uc03}
\begin{tabular}{p{3.2cm}p{10.8cm}}
\toprule
\textbf{Trường} & \textbf{Nội dung} \\
\midrule
Mã & UC-03 \\
Tóm tắt & Người dùng gửi một câu tiếng Việt, nhận câu trả lời kèm Điều dẫn. \\
Trước khi làm & Đã đăng nhập. Chỉ mục vector sẵn thì hỏi được; chưa sẵn thì
nhận 503 và được hướng sang trang tra cứu. \\
Sau khi làm & Câu hỏi và câu trả lời nằm trong hội thoại, kèm JSON trích dẫn
và vết các bước. \\
Luồng chính & 1. Gửi câu. 2. Máy tạo hội thoại nếu chưa có, lấy vài lượt cũ.
3. Ghi câu hỏi. 4. Chạy bảy bước, mỗi bước một sự kiện SSE. 5. Ghi câu trả
lời khi có sự kiện kết quả. 6. Trình duyệt đổi URL sang \texttt{/chat/\{id\}}. \\
Nhánh & Câu không phải pháp luật: thoát ở bước ý định, không tốn truy hồi.
Không tìm ra Điều: nói không có căn cứ, không bịa. \\
Lỗi & Mất kết nối: câu hỏi đã lưu, gửi lại được. Hết hạn mức mô hình: báo lỗi
rõ, không xóa hội thoại. Quá số lần gọi/phút: 429. \\
\bottomrule
\end{tabular}
\end{table}

\begin{table}[htp]
\centering
\caption{Đặc tả UC-07 --- Soát xét hợp đồng.}
\label{tab:uc07}
\begin{tabular}{p{3.2cm}p{10.8cm}}
\toprule
\textbf{Trường} & \textbf{Nội dung} \\
\midrule
Mã & UC-07 \\
Tóm tắt & Tải file, chạy soát xét nền, đọc bảng rủi ro. \\
Trước khi làm & Đã gắn công ty. File \texttt{.pdf}, \texttt{.docx} hoặc
\texttt{.txt}, $\le 10$~MB. \\
Sau khi làm & Tài liệu \texttt{done} hoặc \texttt{failed}. Có thể chạy lại. \\
Luồng chính & 1. Tải, trích chữ, trạng thái \texttt{ready}. 2. Bấm soát xét,
trả 202 ngay. 3. Nền: cắt điều khoản, mỗi khoản đi truy hồi, chấm rủi ro. 4.
Trang hỏi lại mỗi 2,5 giây. \\
Nhánh & PDF ảnh (không có lớp chữ): fail lúc tải, không lúc soát xét, để người
dùng biết ngay file không trích được chữ. Quá nhiều điều khoản: cắt theo ngưỡng
cấu hình, ghi trong tóm tắt. \\
Lỗi & Sai định dạng: 400. Mất mô hình giữa chừng: \texttt{failed}, bấm thử
lại. \\
\bottomrule
\end{tabular}
\end{table}

\subsection{Đặc tả bổ sung UC-05 và UC-09}

\begin{table}[htp]
\centering
\caption{Đặc tả UC-05 --- Tra cứu và đọc văn bản.}
\label{tab:uc05}
\begin{tabular}{p{3.2cm}p{10.8cm}}
\toprule
\textbf{Trường} & \textbf{Nội dung} \\
\midrule
Mã & UC-05 \\
Tóm tắt & Duyệt cây Chương--Điều, tìm chữ (kể cả không dấu), mở viện dẫn chéo. \\
Trước khi làm & Đã đăng nhập. Corpus đã nạp Postgres (không cần Chroma). \\
Sau khi làm & Người dùng đọc được nguyên văn Điều; có thể chuyển sang hỏi đáp. \\
Luồng chính & 1. Mở \texttt{/laws}. 2. Lọc lĩnh vực hoặc chọn số hiệu. 3. Xem cây.
4. Mở Điều. 5. (Tuỳ chọn) bấm viện dẫn hoặc ``Hỏi AI về Điều này''. \\
Nhánh tìm & Gõ ô tìm: ưu tiên số Điều; rồi FTS; rồi trigram nếu không hit. \\
Lỗi & Số hiệu sai đường dẫn: 404. Query rỗng: không gọi API. \\
\bottomrule
\end{tabular}
\end{table}

\begin{table}[htp]
\centering
\caption{Đặc tả UC-09 --- Lịch tuân thủ.}
\label{tab:uc09}
\begin{tabular}{p{3.2cm}p{10.8cm}}
\toprule
\textbf{Trường} & \textbf{Nội dung} \\
\midrule
Mã & UC-09 \\
Tóm tắt & Xem, đánh dấu, sinh lại các mốc nghĩa vụ theo hồ sơ công ty. \\
Trước khi làm & User thuộc một organization có \texttt{vat\_period}, số lao động. \\
Sau khi làm & Task phản ánh đúng rule áp dụng; việc đã tick không bị xóa khi sinh lại. \\
Luồng chính & 1. Đăng ký/cập nhật hồ sơ. 2. Hệ thống sinh cửa sổ 3 tháng trước +
12 tháng tới. 3. User xem theo tháng, đánh dấu xong/bỏ qua. 4. Đổi hồ sơ thì
bấm sinh lại. \\
Nhánh & Công ty $< 10$ lao động: không có rule ATVSLĐ. GTGT tháng/quý loại trừ lẫn nhau. \\
Lỗi & Chưa gắn công ty: API từ chối. Rule thiếu legal\_refs: không được seed. \\
\bottomrule
\end{tabular}
\end{table}

\section{Bảy bước hỏi đáp}

\label{sec:pipeline}

Đồ thị LangGraph gồm bảy nút; tên trong mã nguồn giữ tiếng Anh để khớp mã hiện
thực. Thứ tự chạy cố định, trừ hai nhánh thoát.

\begin{enumerate}
    \item \texttt{analyze\_intent}. Mô hình chỉ được trả \texttt{NEXT} hoặc
    \texttt{SKIP}. Câu chào, hỏi giờ, hỏi ``ăn gì'' dừng tại đây, không tốn
    nhúng. Khi chạy hàng loạt câu hỏi (trang quản trị), nút này luôn \texttt{NEXT}
    để mọi câu đều đi truy hồi.
    \item \texttt{prepare\_retrieval\_query}. Cấu hình sản phẩm: HyDE bật, viết
    lại truy vấn nhiều biến thể tắt. Bật đồng thời cả hai làm tăng độ trễ và hạn mức, trong khi HyDE đã đủ để kéo Điều đúng giọng luật.
    Nếu gọi HyDE lỗi (hết hạn mức, timeout), nút này nuốt lỗi và tìm bằng câu
    gốc --- không được phép làm hỏng cả lượt hỏi chỉ vì một đoạn giả.
    \item \texttt{retrieve}. BM25 và Chroma chạy song song, gộp RRF,
    $k_{\mathrm{RRF}}=60$, trọng số dense/BM25 $= 3/1$, lấy top $k=8$ Điều.
    \item \texttt{rerank}. Trong file cấu hình đang tắt vì Gemini không có cổng
    rerank riêng; nút vẫn nằm trên đồ thị để bật được khi tự host mô hình rerank.
    \item \texttt{llm\_filter}. Tối đa 3 Điều lọc song song, giữ tối thiểu 2
    (tránh lọc sạch rồi không còn gì để viết). Mỗi ứng viên được hỏi ``có dính
    câu hỏi này không'', chỉ được trả có hoặc không.
    \item \texttt{generate\_answer}. Nhiệt độ 0, tối đa 4096 token, stream từng
    token xuống SSE. Prompt hệ thống siết: chỉ kết luận từ căn cứ đang có; thiếu
    căn cứ thì nói thiếu, không bịa mức phạt.
    \item \texttt{format\_submission}. Regex lấy các cụm ``Điều~$n$'' trong câu
    trả lời, đối chiếu với danh sách id vừa truy hồi. Số không có trong danh sách bị
    gỡ khỏi mảng citation gửi xuống giao diện. Nội dung câu trả lời được giữ nguyên để tránh cắt câu; người dùng chỉ không bấm được vào nguồn không thuộc tập truy hồi.
\end{enumerate}

\begin{figure}[htp]
\centering
\resizebox{0.98\linewidth}{!}{\hinhRag}
\caption{Bảy nút LangGraph: nhánh SKIP thoát sớm; rerank tắt trên sản phẩm Gemini.}
\label{fig:hoat-dong-rag}
\end{figure}

\section{Corpus}

\label{sec:corpus}

Kho nằm trong \texttt{data/base\_data.json}: \soDieu{} bản ghi. Tệp kê khai
\texttt{corpus/law\_manifest.json} gắn lĩnh vực cho từng số hiệu, dùng để đổ
bảng \texttt{laws} mỗi lần backend lên.

\subsection{Nguồn ngữ liệu}

Đồ án không gửi request hàng loạt lên cổng văn bản quy phạm. Cả vbpl.vn và
thuvienphapluat.vn đều chặn truy cập tự động. Nội dung được lấy từ bộ dữ liệu
công khai \path{pdt590/vietnamese-legal-documents} trên Hugging
Face~\cite{pdt590}: bản sao dạng markdown của thư viện pháp luật, khoảng 518
nghìn văn bản. Script \path{corpus/build_corpus.py} chỉ giữ các số hiệu có
trong manifest, đọc metadata và mười một phần parquet, kiểm tra đúng văn bản
(từ chối nếu số Điều quá ít --- đã gặp trường hợp mirror trả nhầm văn bản),
rồi cắt thành từng Điều và ghi JSON.

Hai hệ quả. Thứ nhất, từng Điều vẫn truy được về số hiệu và URL gốc (lưu trong
báo cáo dựng kho). Thứ hai, chất lượng kho phụ thuộc bản sao: script không thay
luật sư đối chiếu hiệu lực. Việc chọn bản còn hiệu lực khi hai số hiệu cùng tên
khác năm chưa được tự động hóa.

\begin{table}[htp]
\centering
\caption{Đếm nhanh theo loại và theo lĩnh vực (một văn bản vào đúng một lĩnh vực).}
\label{tab:corpus}
\begin{tabular}{p{3.4cm}r p{4.8cm}r}
\toprule
\textbf{Loại} & \textbf{Số} & \textbf{Lĩnh vực} & \textbf{Số} \\
\midrule
Luật & 19 & Thuế & 12 \\
Nghị định & 14 & Lao động & 4 \\
Thông tư & 3 & Sở hữu trí tuệ & 4 \\
Bộ luật & 2 & Doanh nghiệp & 3 \\
 &  & Dân sự -- hợp đồng & 3 \\
 &  & Kế toán -- tài chính & 3 \\
 &  & Còn lại (BHXH, ATVSLĐ, NTD, đất đai, hỗ trợ DNNVV) & 9 \\
\midrule
\textbf{Cộng} & \textbf{38} & \textbf{Cộng} & \textbf{38} \\
\bottomrule
\end{tabular}
\end{table}

Thuế nhiều nhất vì lịch tuân thủ và câu hỏi hàng ngày của kế toán tập trung ở
đó. Có cả luật cũ và luật mới cùng tên (ví dụ Luật GTGT 13/2008 và 48/2024):
hệ thống giữ cả hai số hiệu vì hợp đồng thực tế vẫn có thể dẫn bản hiệu lực tại thời điểm ký. Việc chọn tự động bản còn hiệu lực chưa được hiện thực và được ghi nhận là hạn chế.

\label{sec:chunking}
Cắt theo Điều đã nói ở Chương~2. Trường \texttt{extra} giữ viện dẫn. Bảng
\texttt{legal\_knowledge\_records} giữ khóa \texttt{id} nguyên từ file JSON,
không tự tăng, để Chroma và Postgres chỉ cùng một mã.

\section{Mô hình dữ liệu}

\begin{figure}[htp]
\centering
\hinhLop
\caption{Kiến trúc ba lớp: trình bày, ứng dụng, dữ liệu.}
\label{fig:so-do-lop}
\end{figure}

\begin{table}[htp]
\centering
\caption{Các bảng nghiệp vụ và ghi chú hiện thực.}
\label{tab:bang-du-lieu}
\begin{tabular}{p{4.0cm}p{9.8cm}}
\toprule
\textbf{Bảng} & \textbf{Ghi chú hiện thực} \\
\midrule
\texttt{organizations} & \texttt{vat\_period}, \texttt{employee\_count} quyết
định rule nào được sinh. \texttt{tax\_code} unique. \\
\texttt{users} & Email lower-case. Role text, không dùng ENUM native Postgres
để khỏi \texttt{ALTER TYPE} mỗi lần thêm giá trị. \\
\texttt{conversations} & Tiêu đề cắt từ câu hỏi đầu, tối đa 80 ký tự, cắt tại
khoảng trắng. \\
\texttt{messages} & Cột \texttt{seq} identity. Không sort theo
\texttt{created\_at}. \\
\texttt{documents} & File trên đĩa tên UUID, không lấy tên gốc. \\
\texttt{contract\_reviews}, \texttt{contract\_findings} & Một tài liệu nhiều
lượt soát xét, không ghi đè. \\
\texttt{compliance\_rules} / \texttt{tasks} & Unique (công ty, rule, kỳ) để
sinh lại không nhân bản. \\
\texttt{legal\_knowledge\_records} & Generated: \texttt{article\_number},
\texttt{doc\_ref}, \texttt{search\_vector}. \\
\texttt{laws} & Đổ từ manifest + đếm Điều, mỗi lần startup. \\
\texttt{audit\_logs} & Đăng ký, đăng xuất, đổi mật khẩu, tải file, reindex. \\
\bottomrule
\end{tabular}
\end{table}

\begin{figure}[p]
\centering
\resizebox{\textwidth}{!}{\hinhCsdl}
\caption{Lược đồ quan hệ PostgreSQL (12 bảng). Hội thoại thuộc user; hợp đồng
và lịch thuộc organization. Corpus cắt theo Điều; \texttt{laws} là catalog đếm
từ cùng \texttt{law\_id}.}
\label{fig:luoc-do-csdl}
\end{figure}

\subsection{Một số quyết định trên lược đồ dữ liệu}

\textbf{Thứ tự tin nhắn.} Câu hỏi và câu trả lời được ghi trong cùng một giao
dịch HTTP. Hàm \texttt{now()} của PostgreSQL trả về thời điểm bắt đầu giao dịch,
nên hai dòng có thể trùng \texttt{created\_at} và thứ tự khi đọc lại không ổn
định. Lược đồ dùng cột \texttt{seq} kiểu identity, API sắp xếp theo
\texttt{seq}.

\textbf{Mã lỗi khi truy cập chéo hội thoại.} Trả 403 khi tài nguyên tồn tại nhưng
không thuộc người dùng sẽ tiết lộ sự tồn tại của định danh. Hệ thống trả 404,
cùng hành vi với trường hợp không có bản ghi.

\textbf{Bộ nhớ hội thoại của LangGraph.} \texttt{InMemorySaver} lưu trạng thái
trên RAM, mất khi khởi động lại và không chia sẻ giữa nhiều worker. Hệ thống nạp
$N$ lượt gần nhất từ PostgreSQL vào state;
\texttt{short\_memory.max\_turns = 6}.

\subsection{Chi tiết từng bảng và khóa ngoại}

Mục này mô tả bảng PostgreSQL phía sản phẩm (không kể bảng Alembic version).
Corpus tra cứu dùng \texttt{legal\_knowledge\_records}; catalog có thể có bảng
\texttt{laws} tuỳ migration. Vector không nêu ở đây.

\subsection{organizations và users}

\texttt{organizations}: tên, mã số thuế unique, số lao động, kỳ GTGT, tuỳ chọn
doanh thu. Xóa tổ chức cascade tài liệu và task; user bị \texttt{SET NULL} khỏi
tổ chức để không mất tài khoản.

\texttt{users}: email unique (lower-case khi ghi), \texttt{password\_hash},
\texttt{role} kiểu text, \texttt{is\_active}. Khóa user không xóa hội thoại ngay
--- tuỳ chính sách admin. UUID làm khóa chính, không lộ số tăng dần.

\subsection{conversations và messages}

Hội thoại thuộc đúng một user, có tiêu đề, cờ lưu trữ, đếm tin nhắn.
\texttt{messages} có \texttt{seq} Identity: cùng giao dịch, \texttt{created\_at}
của Postgres có thể trùng, sắp theo thời gian làm đảo câu hỏi/đáp. JSONB lưu
citation và vết pipeline cho tin assistant. Xóa hội thoại cascade tin nhắn.

\subsection{legal\_knowledge\_records}

Mỗi hàng một Điều. Khóa \texttt{id} lấy từ corpus, không autoincrement, để id
ổn định giữa lần nạp. Cột generated: \texttt{article\_number} (số để sắp 9 trước
10), \texttt{doc\_ref}, \texttt{search\_vector} (tsvector đã unaccent, trọng số
A/B/C theo tiêu đề / tên luật / nội dung). JSONB \texttt{extra}: viện dẫn chéo.
Chỉ mục: (\texttt{law\_id}, \texttt{article}), GIN tsvector, GIN trigram tiêu đề
và tên luật.

\subsection{compliance\_rules và compliance\_tasks}

Rule: mã unique, tần suất, công thức hạn, \texttt{legal\_refs},
\texttt{applies\_to}, \texttt{is\_active}. Task: tổ chức, rule, nhãn kỳ, ngày
đến hạn, trạng thái, unique (tổ chức, rule, kỳ). Index theo hạn để lịch tháng
không quét full.

\subsection{documents, contract\_reviews, contract\_findings}

Document: file trên volume, MIME, kích thước, chữ trích, trạng thái. Review: điểm
rủi ro 0--100, tóm tắt, số điều khoản. Finding: vị trí, mức, gợi ý, căn cứ. Xóa
tài liệu cascade review và finding. Đường file không phải URL công khai; tải về
đi qua API đã xác thực.

\subsection{Nhật ký audit (nếu có trong schema)}

Bảng audit ghi hành động admin (reindex, khóa user) khi module được bật, phục vụ
truy vết, không phục vụ RAG. Không ghi mật khẩu hay token.

\begin{table}[htp]
\centering
\caption{Bảng PostgreSQL chính và khóa ngoại then chốt.}
\label{tab:bang-pg}
\begin{tabular}{p{4.4cm}p{9.3cm}}
\toprule
\textbf{Bảng} & \textbf{Liên kết / ràng buộc} \\
\midrule
\texttt{users} & \texttt{organization\_id} $\rightarrow$ organizations, SET NULL \\
\texttt{conversations} & \texttt{user\_id} CASCADE \\
\texttt{messages} & \texttt{conversation\_id} CASCADE; \texttt{seq} Identity \\
\texttt{legal\_knowledge\_records} & PK id từ corpus; GIN search \\
\texttt{compliance\_tasks} & unique (org, rule, period) \\
\texttt{documents} & org CASCADE; uploaded\_by SET NULL \\
\texttt{contract\_findings} & review CASCADE \\
\bottomrule
\end{tabular}
\end{table}


\section{Tìm chữ trên tiếng Việt}

Hàm \texttt{unaccent} mặc định là \texttt{STABLE}, không dùng trong generated
column. Hàm \texttt{immutable\_unaccent} được bọc trong \texttt{init.sql}, chạy bằng
superuser lúc tạo cụm, Alembic không đụng tới (user ứng dụng không đủ quyền
\texttt{CREATE EXTENSION}).

Ba tầng tìm, theo thứ tự:

\begin{enumerate}
    \item Câu có dạng ``Điều 12'' thì lọc thẳng \texttt{article\_number}.
    \item \texttt{tsquery} kiểu prefix (\texttt{doanh:*}) trên
    \texttt{search\_vector}, xếp \texttt{ts\_rank\_cd}.
    \item Nếu không ra dòng nào, \texttt{pg\_trgm} với ngưỡng 0,3 trên tiêu đề
    Điều và tên văn bản --- cứu gõ sai.
\end{enumerate}

Hàm \texttt{ts\_headline} chạy trên cột đã bỏ dấu, trả về đoạn không dấu nên
không dùng. Backend gửi list token; React bôi \texttt{<mark>} trên chữ gốc
bằng so khớp bỏ dấu phía client.

\section{Luồng thời gian thực và lưu hội thoại}

\begin{figure}[htp]
\centering
\hinhChatSeq
\caption{Luồng hỏi đáp: ghi câu hỏi trước, stream SSE, ghi câu trả lời ở session CSDL riêng.}
\label{fig:tuan-tu-chat}
\end{figure}

Ba loại sự kiện: \texttt{status}, \texttt{token}, \texttt{result}. Có heartbeat
khi một bước kéo dài, để proxy khỏi cắt kết nối. Client không dùng
\texttt{EventSource}.

\section{Hợp đồng}

Cắt điều khoản bằng regex tiêu đề ``Điều n'', ``Article n'', hoặc dòng số kiểu
``5. Thanh toán''. Ưu tiên đuôi file vì trình duyệt hay khai
\texttt{application/octet-stream}. Khoản dưới 80 ký tự ghép vào khoản trước ---
thường là tiêu đề mục. Khoản dài quá 6000 ký tự bị cắt để khỏi vượt ngữ cảnh.
Tối đa 60 khoản một lượt; hợp đồng mẫu DNNVV dùng khi thử nghiệm nằm dưới ngưỡng
này. Mỗi khoản gọi lại cùng hàm truy hồi của hỏi đáp, lấy 4 Điều, rồi mô hình
gán mức \texttt{cao} / \texttt{trung bình} / \texttt{thấp} / \texttt{thông tin}.
Điểm tổng 0--100 từ trọng số mức. PDF không có lớp chữ (scan) fail lúc tải, có
câu tiếng Việt bảo cần OCR --- không đợi đến lúc bấm soát xét mới báo.

\begin{figure}[htp]
\centering
\resizebox{0.95\linewidth}{!}{\hinhHopDongSeq}
\caption{Luồng soát xét hợp đồng: trả 202 ngay, xử lý nền, UI poll trạng thái.}
\label{fig:tuan-tu-hopdong}
\end{figure}

\begin{figure}[htp]
\centering
\hinhTrangThai
\caption{Máy trạng thái tài liệu: pending $\rightarrow$ ready $\rightarrow$ processing $\rightarrow$ done/failed.}
\label{fig:trang-thai-tai-lieu}
\end{figure}

Lúc viết frontend, Axios gắn sẵn \texttt{Content-Type: application/json}, gửi
\texttt{FormData} làm hỏng biên multipart. Interceptor gỡ header đó khi body là
\texttt{FormData}. Đây là lỗi chỉ lộ khi bấm tải file trên trình duyệt; gọi API
thuần JSON thì không dính.

\section{Lịch tuân thủ}

\soRule{} rule nằm trong mã nguồn Python, được seed mỗi lần backend khởi động. Đồ án không
nhét hạn chót vào file cấu hình YAML: rule gắn \texttt{legal\_refs} (số hiệu +
tên văn bản + số Điều). Nếu một hạn chót đổi vì luật mới, người sửa phải đổi
cả refs, không được chỉ sửa ngày.

\begin{table}[htp]
\centering
\caption{Mười hai rule tuân thủ và điều kiện hồ sơ để rule được sinh.}
\label{tab:compliance-rules}
\begin{tabular}{p{4.4cm}p{3.0cm}p{6.4cm}}
\toprule
\textbf{Mã} & \textbf{Kỳ} & \textbf{Điều kiện / hạn chót} \\
\midrule
\texttt{VAT\_MONTHLY} & Tháng & Chỉ khi khai GTGT tháng; ngày 20 tháng sau \\
\texttt{VAT\_QUARTERLY} & Quý & Chỉ khi khai quý; cuối tháng đầu quý sau \\
\texttt{CIT\_PROVISIONAL} & Quý & Mọi công ty; ngày 30 tháng đầu quý sau \\
\texttt{CIT\_FINALIZATION} & Năm & Cuối tháng 3 \\
\texttt{PIT\_FINALIZATION} & Năm & Có từ 1 lao động; cuối tháng 3 \\
\texttt{BUSINESS\_LICENSE\_FEE} & Năm & 30/01 \\
\texttt{SOCIAL\_INSURANCE\_MONTHLY} & Tháng & Có lao động; cuối tháng \\
\texttt{LABOUR\_REPORT\_H1} / \texttt{H2} & Năm & 05/06 và 05/12 \\
\texttt{LABOUR\_SAFETY\_REPORT} & Năm & Từ 10 lao động; 10/01 năm sau \\
\texttt{FINANCIAL\_STATEMENT} & Năm & Cuối tháng 3 \\
\texttt{INVOICE\_USAGE\_CHECK} & Quý & Rà hóa đơn điện tử, ngày 15 tháng sau quý \\
\bottomrule
\end{tabular}
\end{table}

Sinh cửa sổ 3 tháng trước (để thấy việc đã quá hạn) và 12 tháng tới. Khóa duy
nhất (công ty, rule, kỳ) nên bấm ``Sinh lại lịch'' sau khi đổi số lao động không
xóa việc đã đánh dấu xong. Công ty 12 người, khai quý: khi đăng ký hệ thống sinh 39
task --- có GTGT quý, không có GTGT tháng, có báo cáo an toàn lao động vì đủ
ngưỡng 10 người. Công ty 3 người khai tháng thì đảo lại: có rule tháng, không
có báo cáo an toàn. Test backend bắt đúng hai nhánh này.

Hạn chế: hạn chót lấy từ Điều trong corpus tại thời điểm thu thập.
Luật sửa sau đó mà chưa nhập văn bản mới thì lịch vẫn nhắc ngày cũ. Hệ thống
không tự đọc công báo.

\section{Đăng nhập, phiên làm việc, bảo mật}

Mật khẩu băm argon2id (passlib). JWT thuật toán HS256, hai loại token, payload
có \texttt{sub}, \texttt{type}, \texttt{iat}, \texttt{exp}, \texttt{jti}. Access
sống 15 phút, refresh 7 ngày. Gọi API bằng refresh token bị từ chối vì
\texttt{type} không khớp. Interceptor phải phân biệt hai loại token;
đúng ra lỗi đó.

Tài khoản đăng ký đầu tiên, khi bảng \texttt{users} còn trống, được gắn
\texttt{admin}. Cờ \texttt{bootstrap\_first\_admin} bật mặc định. Không có tài
khoản ``admin/admin'' cứng trong image, nhằm tránh mật khẩu mẫu nằm trong
Docker Hub sau này.

Phía trình duyệt, Axios gắn sẵn \texttt{Authorization}. Gặp 401, một lần
refresh được chia cho mọi request đang chờ (một \texttt{Promise} chung), tránh
gọi \texttt{/auth/refresh} song song rồi làm hỏng phiên. Endpoint auth không đi
vào vòng refresh, kẻo login sai mật khẩu cũng bị ``làm mới token''. Token nằm
\texttt{localStorage}: phù hợp quy mô đồ án nhưng còn rủi ro XSS; CSP và cookie \texttt{HttpOnly} chưa được hiện thực.

CORS: danh sách origin, \texttt{allow\_credentials=true} nên không được
\texttt{*}. Trên Docker, nginx và API cùng origin nên CORS gần như không chạy;
danh sách chủ yếu phục vụ Vite lúc lập trình (\texttt{localhost:5173}).

Giới hạn tần suất: chat 20 lần/phút, tải file 30 lần/giờ, theo từng user (hoặc
IP nếu chưa đăng nhập). User thường gọi \texttt{/api/v1/admin/*} nhận 403.
Hội thoại người khác: 404, đã nói ở trên.

Bốn vai trò \texttt{owner}, \texttt{accountant}, \texttt{hr}, \texttt{admin}.
Hợp đồng và lịch tuân thủ cách ly theo \texttt{organization\_id}. Hội thoại
cách ly theo \texttt{user\_id} --- kế toán và chủ cùng công ty không đọc được
chat của nhau. Thiết kế này phù hợp câu hỏi pháp lý mang tính riêng (lương, kỷ luật).

\section{Giao diện}

Frontend là một SPA. Người chưa đăng nhập chỉ vào \texttt{/login} và
\texttt{/register}. Người đã đăng nhập bị đẩy khỏi hai trang đó. Route cần phiên
chờ \texttt{restore()} đọc token xong mới quyết định, nếu không sẽ nháy sang
login rồi nhảy lại.

Số hiệu văn bản kiểu \texttt{45/2019/QH14} chứa dấu gạch chéo. React Router mặc
định cắt path theo \texttt{/}, nên định tuyến bắt phần còn lại bằng
\texttt{/laws/:lawId/*}. Backend cũng nhận path, không nhận query một đoạn, vì
cùng lý do.

Trang quản trị \texttt{lazy()} tách bundle: user thường không tải mã corpus và
trang chạy hàng loạt. TanStack Query tắt \texttt{refetchOnWindowFocus} --- corpus
không đổi mỗi lần đổi tab, refetch chỉ tổ làm nhấp nháy cây Điều.

Ảnh chụp bố cục từng trang (hỏi đáp, tra cứu, hợp đồng, lịch) đặt ở Chương~4 khi
đánh giá giao diện, tránh trùng hình trong báo cáo.

Trang hồ sơ cho sửa số lao động và kỳ GTGT; lưu xong lịch được sinh lại. Trang
chạy hàng loạt: tải file JSON danh sách câu, xem tiến độ, tải kết quả JSON.

Giao diện dùng thành phần tự viết để giảm kích thước gói. Nền tối, chữ đủ tương phản để đọc Điều dài. Bố cục chưa tối ưu cho điện thoại; thanh hội thoại trên màn hẹp còn chật --- đây là hạn chế đã ghi nhận.

\section{Nhúng vector, cache BM25 và biến môi trường}

Mỗi lần đổi model nhúng hoặc đổi số chiều, backend so manifest trong thư mục
Chroma với cấu hình hiện tại. Lệch thì nhúng lại. Lô 64 Điều, nghỉ dần khi gặp
429 (đợi $2$ giây, nhân đôi, trần 60 giây, tối đa 6 lần). Gói miễn phí Google
dễ đụng trần nếu nhúng \soDieu{} Điều một mạch; không có backoff thì lần đầu
cài gần như chắc fail giữa chừng.

BM25 tách từ bằng \texttt{underthesea} (fallback regex). Kết quả ghi file cache
theo dấu vân tay tokenizer. \soDieu{} Điều tách lại mỗi lần restart thì startup
kéo dài hàng phút; cơ chế cache được bổ sung sau khi đo thời gian khởi động.

Docker Compose chỉ đọc \texttt{.env} ở thư mục gốc dự án. File \texttt{backend/.env} dùng khi lập triển khai độc lập \textbf{không} được compose nạp. Khóa API phải đặt tại file gốc; hướng dẫn được nêu trong README và Phụ lục A.

Thiếu khóa: healthcheck vẫn xanh, \texttt{/health} trả ổn, đăng nhập và tra cứu
chạy. Hỏi đáp trả 503 tiếng Việt, gợi ý dùng trang luật. Đây là chỗ thầy góp ý
sớm; phiên bản đầu tiên từng \texttt{raise} lúc import embedding client, kéo theo
cả cụm không lên.

\section{Đóng gói ba dịch vụ}

Ba container trên một mạng cầu. Postgres 17, healthcheck \texttt{pg\_isready},
volume riêng, \texttt{init.sql} tạo extension \texttt{unaccent}, \texttt{pg\_trgm}
và hàm \texttt{immutable\_unaccent}. Backend đợi Postgres healthy rồi mới lên,
healthcheck HTTP \texttt{/health}, \texttt{start\_period} 60 giây vì lần đầu
Alembic và nạp catalog luật chậm. Frontend là nginx phục vụ bản Vite đã build,
proxy \texttt{/api} sang backend, \texttt{VITE\_API\_BASE\_URL} để trống để SPA
gọi đường dẫn tương đối --- cùng origin, hết chuyện CORS trên môi trường cài
đặt.

Cổng ra máy host: web 5173, API 8023, Postgres 23432 (không 5432, tránh đụng
Postgres sẵn có trên máy phát triển). Corpus và \texttt{data/} mount read-only.
\texttt{chroma\_db} và \texttt{uploads} là volume, sống sót khi rebuild image.

\section{Tóm tắt hiện thực theo module}

\begin{table}[htp]
\centering
\caption{Ánh xạ module code $\leftrightarrow$ nghiệp vụ.}
\label{tab:module-map}
\begin{tabular}{p{4.2cm}p{9.5cm}}
\toprule
\textbf{Thư mục / router} & \textbf{Việc} \\
\midrule
\texttt{routers/auth} & Đăng ký, login, refresh, đổi mật khẩu \\
\texttt{routers/legal} + \texttt{agents/} & SSE hỏi đáp, bảy nút LangGraph \\
\texttt{routers/laws} + \texttt{legal/search} & Catalog, cây, FTS ba tầng \\
\texttt{routers/documents} + \texttt{contracts/} & Upload, trích chữ, soát xét nền \\
\texttt{routers/compliance} + \texttt{seed} & \soRule{} rule, sinh task \\
\texttt{routers/admin}, \texttt{lab} & Corpus, user, chạy hàng loạt câu hỏi \\
\texttt{frontend/pages/*} & Chat, Laws, Contracts, Compliance, Admin, bộ câu hỏi \\
\bottomrule
\end{tabular}
\end{table}

Các kịch bản hỏi đáp và chạy cụm trình bày ở Chương~4 --- chương này chỉ chốt
thiết kế và quyết định hiện thực.

\section{Thiết kế API, cấu hình và các khối hiện thực}

\subsection{API theo tài nguyên}

\boncau
{REST dưới tiền tố \texttt{/api/v1}: \texttt{/auth}, \texttt{/conversations},
\texttt{/legal/chat/stream}, \texttt{/laws}, \texttt{/documents},
\texttt{/compliance}, \texttt{/admin}, \texttt{/health}.}
{Tách tài nguyên khớp bốn nghiệp vụ; \texttt{v1} để sau này thêm phiên bản
không phá client cũ.}
{GET đọc, POST tạo, PATCH sửa, DELETE xóa. Lỗi JSON \texttt{detail}.
\texttt{/health} không cần token.}
{Mọi request nghiệp vụ. SSE chỉ trên \texttt{/legal/chat/stream}.}

\subsection{Đăng ký và bootstrap quản trị}

\boncau
{Tài khoản đầu (bảng user trống, cờ \texttt{bootstrap\_first\_admin}) nhận
admin; các tài khoản sau là \texttt{owner} nếu kèm tổ chức.}
{Không nhúng cặp \texttt{admin/admin} trong image.}
{Đăng ký có thể kèm MST, số lao động, kỳ GTGT; ngay sau đó sinh lịch.}
{Chỉ lần đăng ký đầu trên CSDL trống. Các lần sau không tự phong admin.}

\subsection{Cấu hình \texttt{config.yaml}}

\boncau
{File YAML mô tả hành vi pipeline và ngưỡng; không chứa bí mật.}
{Đổi HyDE/rewrite/rerank/filter không sửa mã. Bí mật đi biến môi trường.}
{Nhóm \texttt{llm}, \texttt{embeddings}, \texttt{legal\_assistant.*},
\texttt{vector\_store} (hybrid, $k=8$, RRF 60, trọng số 3:1),
\texttt{short\_memory.max\_turns=6}, \texttt{auth}, \texttt{rate\_limit},
\texttt{uploads} (10~MB).}
{Mỗi lần backend khởi động. Đổi chiều embedding: xóa volume Chroma rồi nhúng lại.}

\subsection{Tách chữ và soát xét hợp đồng}

\boncau
{Trích chữ PDF (lớp văn bản), DOCX (XML), TXT (UTF-8); cắt theo tiêu đề điều
khoản; mỗi khoản đi cùng hàm truy hồi rồi LLM gán mức rủi ro.}
{Tái sử dụng retrieval: căn cứ hợp đồng cũng mở được Điều, cùng hàng rào với chat.
Không huấn luyện mô hình phát hiện rủi ro riêng.}
{Regex ``Điều $n$'', ``Article $n$'', dòng số. Gộp khoản $<80$ ký tự; cắt khoản
$>6000$. Tối đa 60 khoản, 3 khoản song song. $k=4$ Điều mỗi khoản. Điểm 0--100
từ trọng số cao/trung/thấp.}
{Trích chữ lúc tải (201/\texttt{ready}). Soát xét lúc user bấm (202, nền). PDF
ảnh: fail lúc tải, không đợi soát xét.}

\subsection{Giao diện từng trang}

\boncau
{SPA: login/register, chat, luật, hợp đồng, lịch, hồ sơ, admin (corpus, user,
chạy hàng loạt).}
{Mỗi nghiệp vụ một nhóm route; citation là nút mở Điều, không phải footnote tĩnh.}
{AppShell ẩn mục admin nếu không phải quản trị. Số hiệu có \texttt{/} dùng splat.}
{Sau đăng nhập. Login/register khi chưa có phiên.}

\subsection{Nhật ký và chiến lược lỗi}

\boncau
{\texttt{audit\_logs} ghi đăng ký, đăng xuất, đổi mật khẩu, tải tệp, reindex.
Thông báo lỗi tiếng Việt, không traceback.}
{Truy vết khi nhầm reindex; người dùng biết thiếu index hay file không có chữ.}
{503 gợi ý tra cứu khi chưa nhúng. Hết hạn mức: giữ câu hỏi đã lưu. Mất SSE:
gửi lại được.}
{Mỗi hành động nhạy cảm (ghi log). Mọi lỗi trả về client (thông báo). Chưa có
UI xem log.}

\section{Danh mục văn bản trong corpus}

\boncau
{Bảng~\ref{tab:danh-muc-38}: đủ \soLuat{} số hiệu đã nạp, mỗi dòng một lĩnh vực
theo manifest.}
{Người đọc đối chiếu được kho thật, không chỉ con số ``38 luật''.}
{Tên rút gọn nếu dài; tên đầy đủ nằm bảng \texttt{laws}.}
{Khi đọc báo cáo và khi admin mở kho. Không phải danh mục VBQPPL toàn quốc.}

\begin{longtable}{p{3.3cm}p{2.0cm}p{2.6cm}p{5.6cm}}
\caption{Danh mục \soLuat{} văn bản quy phạm trong corpus.}
\label{tab:danh-muc-38}\\
\toprule
\textbf{Số hiệu} & \textbf{Loại} & \textbf{Lĩnh vực} & \textbf{Tên} \\
\midrule
\endfirsthead
\caption[]{Danh mục văn bản (tiếp).}\\
\toprule
\textbf{Số hiệu} & \textbf{Loại} & \textbf{Lĩnh vực} & \textbf{Tên} \\
\midrule
\endhead
\texttt{04/2017/QH14} & Luật & Hỗ trợ DNNVV & Luật Hỗ trợ doanh nghiệp nhỏ và vừa \\
\texttt{80/2021/NĐ-CP} & Nghị định & Hỗ trợ DNNVV & Nghị định Quy định chi tiết và hướng dẫn thi hành một số điều của Luật Hỗ trợ doanh ngh... \\
\texttt{59/2020/QH14} & Luật & Doanh nghiệp & Luật Doanh nghiệp \\
\texttt{01/2021/NĐ-CP} & Nghị định & Doanh nghiệp & Nghị định về đăng ký doanh nghiệp \\
\texttt{122/2021/NĐ-CP} & Nghị định & Doanh nghiệp & Nghị định Quy định về xử phạt vi phạm hành chính trong lĩnh vực kế hoạch và đầu tư \\
\texttt{38/2019/QH14} & Luật & Thuế & Luật Quản lý thuế \\
\texttt{108/2025/QH15} & Luật & Thuế & Luật Quản lý thuế \\
\texttt{126/2020/NĐ-CP} & Nghị định & Thuế & Nghị định Quy định chi tiết một số điều của Luật Quản lý thuế \\
\texttt{123/2020/NĐ-CP} & Nghị định & Thuế & Nghị định Quy định về hóa đơn, chứng từ \\
\texttt{78/2021/TT-BTC} & Thông tư & Thuế & Thông tư Hướng dẫn thực hiện một số điều của Luật Quản lý thuế ngày 13 tháng 6 năm 2019... \\
\texttt{132/2020/NĐ-CP} & Nghị định & Thuế & Nghị định Quy định về quản lý thuế đối với doanh nghiệp có giao dịch liên kết \\
\texttt{139/2016/NĐ-CP} & Nghị định & Thuế & Nghị định Quy định về lệ phí môn bài \\
\texttt{13/2008/QH12} & Luật & Thuế & Luật Thuế giá trị gia tăng \\
\texttt{48/2024/QH15} & Luật & Thuế & Luật Thuế giá trị gia tăng \\
\texttt{14/2008/QH12} & Luật & Thuế & Luật Thuế thu nhập doanh nghiệp \\
\texttt{67/2025/QH15} & Luật & Thuế & Luật Thuế thu nhập doanh nghiệp \\
\texttt{04/2007/QH12} & Luật & Thuế & Luật Thuế thu nhập cá nhân \\
\texttt{45/2019/QH14} & Bộ luật & Lao động & Bộ luật Lao động \\
\texttt{145/2020/NĐ-CP} & Nghị định & Lao động & Nghị định Quy định chi tiết và hướng dẫn thi hành một số điều của Bộ luật Lao động về đ... \\
\texttt{12/2022/NĐ-CP} & Nghị định & Lao động & Nghị định Quy định xử phạt vi phạm hành chính trong lĩnh vực lao động, bảo hiểm xã hội,... \\
\texttt{152/2020/NĐ-CP} & Nghị định & Lao động & Nghị định Quy định về người lao động nước ngoài làm việc tại Việt Nam và tuyển dụng, qu... \\
\texttt{41/2024/QH15} & Luật & Bảo hiểm xã hội & Luật Bảo hiểm xã hội \\
\texttt{58/2020/NĐ-CP} & Nghị định & Bảo hiểm xã hội & Nghị định Quy định về mức đóng bảo hiểm xã hội bắt buộc vào Quỹ bảo hiểm tai nạn lao độ... \\
\texttt{84/2015/QH13} & Luật & An toàn lao động & Luật An toàn, vệ sinh lao động \\
\texttt{44/2016/NĐ-CP} & Nghị định & An toàn lao động & Nghị định Quy định chi tiết một số điều của Luật An toàn, vệ sinh lao động về hoạt động... \\
\texttt{50/2005/QH11} & Luật & Sở hữu trí tuệ & Luật Sở hữu trí tuệ \\
\texttt{36/2009/QH12} & Luật & Sở hữu trí tuệ & Luật sửa đổi, bổ sung một số điều của Luật Sở hữu trí tuệ \\
\texttt{07/2022/QH15} & Luật & Sở hữu trí tuệ & Luật sửa đổi, bổ sung một số điều của Luật Sở hữu trí tuệ \\
\texttt{65/2023/NĐ-CP} & Nghị định & Sở hữu trí tuệ & Nghị định Quy định chi tiết một số điều và biện pháp thi hành Luật Sở hữu trí tuệ về sở... \\
\texttt{91/2015/QH13} & Bộ luật & Dân sự - Hợp đồng & Bộ luật Dân sự \\
\texttt{36/2005/QH11} & Luật & Dân sự - Hợp đồng & Luật Thương mại \\
\texttt{54/2010/QH12} & Luật & Dân sự - Hợp đồng & Luật Trọng tài thương mại \\
\texttt{88/2015/QH13} & Luật & Kế toán - Tài chính & Luật Kế toán \\
\texttt{133/2016/TT-BTC} & Thông tư & Kế toán - Tài chính & Thông tư Hướng dẫn Chế độ kế toán doanh nghiệp nhỏ và vừa \\
\texttt{200/2014/TT-BTC} & Thông tư & Kế toán - Tài chính & Thông tư Hướng dẫn Chế độ kế toán doanh nghiệp \\
\texttt{19/2023/QH15} & Luật & Bảo vệ người tiêu dùng & Luật Bảo vệ quyền lợi người tiêu dùng \\
\texttt{98/2020/NĐ-CP} & Nghị định & Bảo vệ người tiêu dùng & Nghị định Quy định xử phạt vi phạm hành chính trong hoạt động thương mại, sản xuất, buô... \\
\texttt{31/2024/QH15} & Luật & Đất đai & Luật Đất đai \\
\bottomrule
\end{longtable}

Danh mục cho thấy thuế chiếm nhiều văn bản nhất, phù hợp lịch tuân thủ và câu
hỏi hàng ngày của kế toán. Có cả luật cũ và luật mới cùng tên (ví dụ quản lý
thuế, thuế GTGT): hệ thống giữ theo số hiệu, chưa tự lọc bản còn hiệu lực.

\section{Mười hai nghĩa vụ tuân thủ và cách sinh lịch}

\boncau
{Bộ \soRule{} bản ghi khai báo trong mã (\texttt{COMPLIANCE\_RULES}): mã, tiêu đề,
tần suất, hạn chót, \texttt{applies\_to}, \texttt{legal\_refs} trỏ Điều trong corpus.}
{DNNVV cần nhắc mốc thuế, BHXH, lao động, BCTC. LLM không được suy hạn nộp ---
dễ bịa ngày.}
{Bộ sinh đọc hồ sơ tổ chức, tạo task cửa sổ 3 tháng trước + 12 tháng tới.
\texttt{INSERT \ldots ON CONFLICT DO NOTHING} theo (công ty, rule, kỳ).}
{Lúc đăng ký/cập nhật hồ sơ và khi user bấm sinh lại. Không chạy khi chưa gắn
tổ chức. Luật đổi hạn: sửa rule trong mã, không chờ mô hình nhớ.}

Mỗi mục dưới đây: nghĩa vụ là gì, công thức hạn trong mã, và khi nào rule được
lọc vào hồ sơ.

\subsection{VAT\_MONTHLY --- khai nộp GTGT theo tháng}

Áp dụng khi hồ sơ khai \texttt{vat\_period = monthly}. Hạn: ngày 20 của tháng
liền sau tháng phát sinh, theo hướng Điều 44 Luật Quản lý thuế (hồ sơ khai) và
Điều 55 (thời hạn nộp tiền) số 38/2019/QH14. \texttt{due\_day} $= 20$,
\texttt{due\_month\_offset} $= 1$.

Doanh nghiệp nhỏ thường được chọn quý; rule tháng tồn tại vì một số đơn vị vẫn
kê khai tháng, và vì demo cần thấy khác biệt khi đổi hồ sơ.

\subsection{VAT\_QUARTERLY --- khai nộp GTGT theo quý}

Áp dụng \texttt{vat\_period = quarterly}. Hạn được mã hoá ngày cuối tháng đầu
quý sau (\texttt{due\_day} $= 31$ kết hợp offset một tháng, bộ sinh kẹp theo số
ngày thật của tháng). Căn cứ: Điều 44 Luật Quản lý thuế và Điều 8 Nghị định
126/2020/NĐ-CP. Đây là rule người dùng DNNVV gặp nhiều nhất trên lịch.

\subsection{CIT\_PROVISIONAL --- tạm nộp thuế TNDN theo quý}

Mọi tổ chức trong phạm vi đồ án. Hạn ngày 30 tháng đầu quý sau. Rule nhấn:
không phát sinh hồ sơ khai tạm tính quý, chỉ nghĩa vụ tiền. Căn cứ Điều 55 Luật
Quản lý thuế và Điều 8 Nghị định 126/2020/NĐ-CP. Nhầm lẫn thường gặp là coi tạm
nộp như một tờ khai --- mô tả rule ghi rõ để UI không gọi là ``khai quý TNDN''.

\subsection{CIT\_FINALIZATION --- quyết toán TNDN năm}

Hạn ngày cuối tháng thứ ba sau khi kết thúc năm dương lịch
(\texttt{due\_fixed\_month} $= 3$, \texttt{due\_day} $= 31$). Căn cứ Điều 44 Luật
Quản lý thuế và Điều 12 Luật Thuế TNDN 14/2008/QH12 (kho còn số hiệu này; luật
mới cùng chủ đề được nạp riêng, rule chưa tự chuyển). Năm tài chính khác dương
lịch chưa tham số hoá --- hạn chế phạm vi.

\subsection{PIT\_FINALIZATION --- quyết toán TNCN năm}

Chỉ khi \texttt{min\_employees} $\ge 1$. Tổ chức trả thu nhập quyết toán thay
người lao động, hạn cùng cửa sổ tháng thứ ba. Căn cứ Điều 44 Luật Quản lý thuế
và Điều 24 Luật Thuế TNCN 04/2007/QH12. Công ty không có lao động không thấy
task này.

\subsection{BUSINESS\_LICENSE\_FEE --- lệ phí môn bài}

Hạn 30 tháng 01 hằng năm. Căn cứ Điều 5 Nghị định 139/2016/NĐ-CP và Điều 10
Nghị định 126/2020/NĐ-CP. Rule không tính số tiền theo vốn điều lệ --- đồ án
nhắc hạn, không thay phần mềm thuế.

\subsection{SOCIAL\_INSURANCE\_MONTHLY --- đóng BHXH hằng tháng}

Khi có ít nhất một lao động. Hạn ngày cuối tháng đang diễn ra
(\texttt{due\_month\_offset} $= 0$, \texttt{due\_day} $= 31$ kẹp cuối tháng).
Căn cứ Điều 34 Luật BHXH 41/2024/QH15 và Điều 4 Nghị định 58/2020/NĐ-CP (quỹ
tai nạn lao động). Tỷ lệ đóng không nằm trong rule; người dùng hỏi tỷ lệ thì
đi pipeline RAG, không đi lịch.

\subsection{LABOUR\_REPORT\_H1 và H2 --- báo cáo tình hình lao động}

Hai rule cùng căn cứ Điều 12 Bộ luật Lao động 45/2019/QH14 và Điều 4 Nghị định
145/2020/NĐ-CP, khác tháng: trước 05/06 và trước 05/12. Lọc
\texttt{min\_employees} $\ge 1$. Tách H1/H2 thay vì một rule ``hai lần/năm''
để task hiện thành hai mốc trên lịch tháng, dễ đánh dấu độc lập.

\subsection{LABOUR\_SAFETY\_REPORT --- báo cáo an toàn vệ sinh lao động}

Khi \texttt{min\_employees} $\ge 10$. Hạn 10 tháng 01 năm sau. Căn cứ Điều 81
Luật An toàn, vệ sinh lao động 84/2015/QH13. Ngưỡng 10 lao động là điều kiện
lọc trong mã, phản ánh hướng tập trung DNNVV có quy mô đủ để nghĩa vụ báo cáo
này trở nên thường gặp; khi văn bản đổi ngưỡng phải sửa \texttt{applies\_to}.

\subsection{FINANCIAL\_STATEMENT --- báo cáo tài chính năm}

Mọi tổ chức. Hạn mã hoá cuối tháng 3 (90 ngày sau 31/12 với năm dương lịch).
Căn cứ Điều 29 Luật Kế toán 88/2015/QH13 và Điều 80 Thông tư 133/2016/TT-BTC
(chế độ kế toán DNNVV). Nộp cơ quan thuế và thống kê được nêu trong mô tả;
hệ thống không tạo hai task tách cơ quan.

\subsection{INVOICE\_USAGE\_CHECK --- rà soát hóa đơn điện tử}

Rule quý, hạn ngày 15 tháng đầu quý sau: đây là \textbf{mốc nội bộ} nhắc rà
soát việc lập, gửi, lưu hóa đơn theo Nghị định 123/2020/NĐ-CP Điều 9 và Thông
tư 78/2021/TT-BTC Điều 6, không phải một tờ khai riêng có số hiệu thống nhất
như GTGT. Mục đích: kéo người dùng từ lịch sang tra cứu/hỏi đáp khi sắp đến
kỳ, chứ không tự đối soát dữ liệu tổng cục thuế.

\subsection{Bảng tổng hợp rule}

\begin{table}[htp]
\centering
\caption{\soRule{} rule tuân thủ: tần suất và điều kiện lọc.}
\label{tab:12-rule}
\begin{tabular}{p{3.6cm}p{2.2cm}p{2.4cm}p{5.2cm}}
\toprule
\textbf{Mã} & \textbf{Tần suất} & \textbf{Lọc hồ sơ} & \textbf{Hướng hạn trong mã} \\
\midrule
VAT\_MONTHLY & Tháng & Kỳ GTGT tháng & Ngày 20 tháng sau \\
VAT\_QUARTERLY & Quý & Kỳ GTGT quý & Cuối tháng đầu quý sau \\
CIT\_PROVISIONAL & Quý & Không & Ngày 30 tháng đầu quý sau \\
CIT\_FINALIZATION & Năm & Không & Cuối tháng 3 \\
PIT\_FINALIZATION & Năm & Có NLĐ & Cuối tháng 3 \\
BUSINESS\_LICENSE\_FEE & Năm & Không & 30 tháng 01 \\
SOCIAL\_INSURANCE\_MONTHLY & Tháng & Có NLĐ & Cuối tháng \\
LABOUR\_REPORT\_H1 & Năm & Có NLĐ & 05 tháng 6 \\
LABOUR\_REPORT\_H2 & Năm & Có NLĐ & 05 tháng 12 \\
LABOUR\_SAFETY\_REPORT & Năm & $\ge$ 10 NLĐ & 10 tháng 01 năm sau \\
FINANCIAL\_STATEMENT & Năm & Không & Cuối tháng 3 \\
INVOICE\_USAGE\_CHECK & Quý & Không & Ngày 15 tháng đầu quý sau \\
\bottomrule
\end{tabular}
\end{table}

Người dùng đổi số lao động hoặc kỳ GTGT rồi bấm sinh lại: task tương lai theo
rule mới xuất hiện; task quá khứ đã đánh dấu xong giữ nguyên. Không dùng LLM để
``đoán'' hạn nộp --- sai một ngày là sai nghiệp vụ.

\section{Quy trình chuẩn hóa ngữ liệu}

\boncau
{Mỗi phần tử JSON là một Điều: số hiệu, chương, số Điều, tiêu đề, nội dung,
\texttt{extra}. Manifest liệt kê văn bản được phép nạp.}
{Cắt theo Điều giữ cách dẫn chiếu của nghề. Cắt chồng theo 512 token phá số
Điều. Khóa \texttt{id} ổn định để Postgres và Chroma cùng một mã.}
{Script nạp kiểm trường bắt buộc và số hiệu khớp manifest. Điều dài vẫn một
vector; Điều ngắn (định nghĩa) vẫn nhúng vì người dùng hay hỏi định nghĩa.}
{Khi dựng CSDL hoặc khi quản trị nhập file. Script không đánh dấu ``còn hiệu
lực'' --- đó là bước biên tập người.}

Thứ tự làm việc khi máy trắng: chọn số hiệu trong manifest, dựng
\texttt{base\_data.json} từ bộ Hugging Face, nạp Postgres, rồi mới nhúng
vector. Đảo thứ tự (nhúng trước khi có hàng chữ) cho chỉ mục rỗng hoặc lệch id.
Nạp chữ và nhúng tách nhau: sau bước Postgres, trang tra cứu đã dùng được, chưa
cần khóa mô hình.

